% 398_Dichte_Selbstwechselwirkung_De_ch.tex
% Johann Pascher, 18. Nov. 2026
% |psi|^2 als Energiedichte und der Einwand der Selbstwechselwirkung

\providecommand{\BM}{\textbf{[B]}}
\providecommand{\KM}{\textbf{[K]}}
\providecommand{\SM}{\textbf{[S]}}
\providecommand{\XM}{\textbf{[X]}}
\providecommand{\QM}{\textbf{[Q]}}
\providecommand{\EM}{\textbf{[E]}}

\chapter*{$|\psi|^2$ als Energiedichte und der Einwand der Selbstwechselwirkung}
\addcontentsline{toc}{chapter}{$|\psi|^2$ als Energiedichte und der Einwand der Selbstwechselwirkung}

\begin{abstract}
In der FFGFT ist $|\psi|^2$ keine Wahrscheinlichkeit, sondern eine Energiedichte; die Born-Werte
erscheinen als geometrische Quotienten (Dok.~020, 129, 230). Diese Lesart hat einen berühmten
Vorläufer, Schrödingers Ladungsdichte $-e|\psi|^2$ von 1926, und einen ebenso bekannten Einwand: Wäre
die Dichte real, müssten ihre Teile aufeinander wirken, und diese Selbstwechselwirkung ist nicht
beobachtet. Das Dokument stellt Schrödingers Lesart und Gaos Einwand mit wörtlichen Zitaten dar und
rechnet ihn nach. Für das Wasserstoffatom ergäbe die Coulomb-Selbstenergie der 1s-Dichte 17{,}0\,eV,
und der 1S--2S-Abstand verschöbe sich um 12{,}9\,eV, mehr als der Übergang selbst, der auf etwa
$4\times10^{-15}$ genau gemessen ist. Eine im Ortsraum ausgedehnte, sich selbst abstoßende
Ladungswolke ist damit ausgeschlossen. Der Einwand setzt aber zweierlei voraus: dass die Dichte als
Kontinuum von Teilen im Ortsraum vorliegt und dass jeder Teil eine Quelle für die anderen ist. Die
FFGFT teilt beides nicht. Ihre Dichte ist die Energiedichte der inneren Konfiguration einer Mode
(Dok.~020, 230), eine Mode ist eine Quelle mit einer Zeitskala (Dok.~334), und das Feld folgt der
realisierten Lage (Dok.~020). In der Feldtheorie entspricht das dem normalgeordneten
Wechselwirkungsterm, der im Einteilchensektor verschwindet; das ist exakt gezeigt. Dass die
FFGFT-Lesart diese Eigenschaft aus der Struktur des Trägers selbst liefert, bleibt offen. Als
Prüfstein ergibt sich: Die FFGFT erwartet keine Schrödinger-Newton-Selbstgravitation.
\end{abstract}

% ============================================================
\section*{Statusmarkierungen}
% ============================================================
\begin{center}
\begin{tabular}{@{}lp{11cm}@{}}
\textbf{[B]} & Algebraisch bewiesen. \\[2pt]
\textbf{[K]} & Numerisch bestätigt (Prüfskript, Vergleich mit Messwerten oder Literaturwerten). \\[2pt]
\textbf{[S]} & Setzung, offen oder Einordnung ohne eigene Rechnung. \\[2pt]
\textbf{[X]} & Ausgeschlossen. \\[2pt]
\textbf{[Q]} & Aus der Literatur übernommen, mit Quelle. \\[2pt]
\textbf{[E]} & Etablierte Physik, im Dokument nicht hergeleitet. \\
\end{tabular}
\end{center}
Prüfskript: \texttt{2/python/Dok398\_Skripte/pruef\_398\_dichte\_selbstwechselwirkung.py}.
Wörtliche Zitate stehen in der Originalsprache in Anführungszeichen und sind mit \QM{} markiert.

% ============================================================
\section{Die Lesart der FFGFT}
% ============================================================

Der Korpus legt die Bedeutung von $|\psi|^2$ an mehreren Stellen gleichlautend fest. Die Amplitude
folgt der lokalen Energiedichte der inneren Konfiguration des Zustands, keiner Wahrscheinlichkeit
(Dok.~020, 384); $|\psi|^2$ ist eine Energiedichte, und die Born-Werte erscheinen als geometrische
Quotienten (Dok.~129, 230). In der Zylinderdarstellung von Dok.~230 ist $|\alpha|^2=(1+z)/2$ der
Anteil der Achse oberhalb des Punkts, und $r^2=1-z^2=4|\alpha|^2|\beta|^2$ ist die Energiedichte der
radialen Konfiguration \BM. Die Zahlenwerte sind dieselben wie in der Quantenmechanik; verschieden
ist, was sie bedeuten. Die Messung ist ein deterministischer Polübergang, dessen Ausgang durch $z$ und
eine in der Wechselwirkung mit der Messanordnung festgelegte Größe $\lambda$ bestimmt ist.

Zur Gravitation hält Dok.~020 fest: Das Einzelergebnis wird in der Wechselwirkung mit der
Messanordnung festgelegt, und das Gravitationsfeld folgt der realisierten Lage. Die Masse ist nicht
als Mittel über zwei Lagen verteilt, wie es die semiklassische Gravitation mit dem Erwartungswert
des Energie-Impuls-Tensors als Quelle annimmt; das passt zu dem Versuch von Page und Geilker
\QM~\cite{PageGeilker1981}. Wie eine Masse vor jeder Wechselwirkung gravitativ wirkt, ist dort
ausdrücklich nicht hergeleitet.

Wer $|\psi|^2$ als reale Dichte liest, muss sich einem alten Einwand stellen. Dieses Dokument stellt
ihn dar, rechnet ihn nach und prüft, ob er die Lesart der FFGFT trifft.

% ============================================================
\section{Schrödingers Ladungswolke}
% ============================================================

Schrödinger hat 1926, im dritten und vierten Teil seiner Wellenmechanik, die Größe $-e|\psi|^2$ als
die Ladungsdichte des Elektrons im Raum gedeutet \QM~\cite{Schroedinger1926}: Das Elektron ist in
dieser Lesart kein Punkt, sondern über den ganzen Raum verteilt, mit einer Dichte, die dem Betrag
der Welle im Quadrat folgt. Er zeigte, dass diese Dichte die richtigen Strahlungsfrequenzen und
Intensitäten liefert, und hielt an der Deutung auch in seinem Bericht für die Solvay-Konferenz 1927
fest \QM~\cite{BacciagaluppiValentini2009}. Durchgesetzt hat sich Borns Deutung als
Wahrscheinlichkeitsdichte.

Die Lesart ist nie ganz verschwunden. Die geschützte Messung von Aharonov, Anandan und Vaidman zeigt,
dass sich $|\psi|^2$ an einem einzelnen System messen lässt, nicht nur als Häufigkeit über viele
Läufe \QM~\cite{AharonovVaidman1993,AAV1993}: Wird der Zustand während einer schwachen, langen
Kopplung geschützt, liefert der Zeiger den Mittelwert der Dichte über ein kleines Gebiet. Gao fasst
das so zusammen, dass Masse und Ladung eines einzelnen Systems im Raum verteilt sind, mit den
Dichten $m|\psi|^2$ und $Q|\psi|^2$ \QM~\cite{Gao2011}; die Geschichte von Schrödingers Hypothese hat er
eigens aufgearbeitet \QM~\cite{Gao2017}. Carver Mead hat die Wellenfunktion in seiner
\enquote{Collective Electrodynamics} ausdrücklich als Materiewelle mit Ladungsdichte behandelt
\QM~\cite{Mead2000}, und Cramer und Mead haben den Kollaps mit Schrödingers ursprünglicher
Wellenmechanik als Austausch zwischen Sender und Empfänger nachgerechnet \QM~\cite{CramerMead2020}.

Für die FFGFT ist die geschützte Messung ein Argument, das in ihre Richtung zeigt: Die Dichte ist am
Einzelsystem eine Größe, nicht erst an der Messreihe \SM.

% ============================================================
\section{Der Einwand}
% ============================================================

Gao formuliert den Einwand so: Wenn die Massen- und Ladungsdichte real ist, also die Dichten an
verschiedenen Orten zugleich vorhanden sind, \enquote{there will exist gravitational and
electrostatic self-interactions of the density} \QM~\cite{Gao2011}. Die Schrödinger-Gleichung eines
freien Systems mit Masse $m$ und Ladung $Q$ müsste dann einen Zusatzterm tragen,
\begin{equation}
  i\hbar\,\partial_t\psi=-\frac{\hbar^2}{2m}\nabla^2\psi
  +\left(kQ^2-Gm^2\right)\!\int\!\frac{|\psi(x',t)|^2}{|x-x'|}\,d^3x'\;\psi ,
  \label{eq:gao}
\end{equation}
mit der Coulomb-Konstante $k$ \QM~\cite{Gao2011}. Als Maß der Stärke nennt er nach Salzman für die
Gravitation $\varepsilon^2=(4Gm^2/\hbar c)^2$ und für die Elektrostatik $\varepsilon^2=(4kQ^2/\hbar
c)^2$; für ein freies Elektron ist das elektrostatische Maß etwa $10^{-3}$, das gravitative etwa
$4\times10^{-89}$ \QM~\cite{Gao2011,Salzman2005}. Die elektrostatische Selbstwechselwirkung
müsste also deutlich sichtbar sein, in Interferenzversuchen mit Elektronen ebenso wie in den
Energieniveaus des Wasserstoffs, wo das Selbstpotential von derselben Größenordnung ist wie das
Potential des Kerns. Da beides nicht beobachtet ist, schließt Gao, die Dichte könne
\enquote{not be real but be effective} \QM~\cite{Gao2011}: In jedem Augenblick gebe es nur ein
lokalisiertes Teilchen, und die Dichte entstehe als Zeitmittel einer zufälligen, unstetigen Bewegung
dieses Teilchens.

% ============================================================
\section{Nachrechnung}
% ============================================================

\textbf{Stärkemaße.} Mit $kQ^2=\alpha\hbar c$ ist das elektrostatische Maß $(4\alpha)^2=8{,}5\times10^{-4}$,
das gravitative $(4Gm_e^2/\hbar c)^2=4{,}9\times10^{-89}$ \KM. Gaos Angaben stimmen in der
Größenordnung. Das Verhältnis beider Maße ist rund $10^{85}$ \KM.

\textbf{Wasserstoff.} Für die 1s-Dichte $\rho(r)=e^{-2r}/\pi$ in atomaren Einheiten ist das
Selbstpotential einer Kugelschale geschlossen bekannt, und die Coulomb-Selbstenergie
\begin{equation}
  J_{1s}=\int\!\!\int\frac{\rho(r)\rho(r')}{|r-r'|}\,d^3r\,d^3r'=\tfrac58\ \text{Hartree}=17{,}0\ \text{eV}
\end{equation}
ist größer als die Bindungsenergie von $13{,}6$\,eV \KM. Für die 2s-Dichte ergibt die numerische
Integration $J_{2s}=77/512$\,Hartree$\,=4{,}09$\,eV \KM. Der Zusatzterm in~\eqref{eq:gao} verschöbe die
Niveaus in erster Ordnung um diese Beträge, also den 1S--2S-Abstand um $12{,}9$\,eV. Der Übergang
selbst beträgt $10{,}2$\,eV; die Verschiebung wäre 127\,\% davon \KM. Die 1S--2S-Frequenz ist auf
etwa 10\,Hz bei 2466\,THz gemessen, also auf rund $4\times10^{-15}$ \QM~\cite{Parthey2011}. Eine
Ladungswolke im Ortsraum, deren Teile sich nach~\eqref{eq:gao} gegenseitig abstoßen, ist damit
ausgeschlossen \XM.

\textbf{Gravitation.} Dieselbe Dichte ergäbe eine gravitative Selbstenergie von $4\times10^{-42}$\,eV,
um den Faktor $Gm_e^2/(kQ^2)=2{,}4\times10^{-43}$ kleiner \KM. Für ein Elektron ist der gravitative
Teil also unbeobachtbar. Bei schweren, ausgedehnten Körpern ist das anders; dafür gibt es die
Schrödinger-Newton-Gleichung, die genau diesen Term für die Gravitation behält
\QM~\cite{Diosi1984,Penrose1996}, und Vorschläge, ihn optomechanisch zu messen
\QM~\cite{Yang2013,Grossardt2016}. Großardt et al.\ beschreiben die Gleichung als
\enquote{experimentally testable alternative to quantum gravity, accessible at low energies}
\QM~\cite{Grossardt2016}.

% ============================================================
\section{Was der Einwand voraussetzt}
% ============================================================

Der Einwand ist zwingend, wenn zwei Annahmen gelten.

\textbf{Erstens:} Die Dichte liegt als Kontinuum von Teilen im Ortsraum vor, und diese Teile sind zur
selben Zeit an verschiedenen Orten. Das ist Schrödingers Wolke im wörtlichen Sinn.

\textbf{Zweitens:} Jeder Teil ist eine Quelle eines Feldes, das auf die anderen Teile desselben
Systems wirkt. Das ist die Annahme, die Gleichung~\eqref{eq:gao} hinschreibt: Die eigene Dichte
erscheint als Quelle im eigenen Potential. Für die Gravitation ist das die semiklassische Kopplung
nach Møller und Rosenfeld, mit dem Erwartungswert als Quelle \EM.

Fällt eine der beiden Annahmen, verliert der Einwand seine Kraft. Gao lässt die erste fallen: Die
Dichte wird effektiv, entstanden aus einer zufälligen Bewegung. Das kostet den Determinismus, denn
die Bewegung ist bei ihm ausdrücklich zufällig und unstetig \QM~\cite{Gao2011}. Für die FFGFT kommt
dieser Ausweg nicht in Frage.

% ============================================================
\section{Die Antwort der FFGFT}
% ============================================================

\subsection{Wo die Dichte liegt}

Die Energiedichte der FFGFT ist nicht die Dichte einer Wolke im Ortsraum, sondern die der inneren
Konfiguration einer Mode (Dok.~020). In der Zylinderdarstellung ist sie $r^2=4|\alpha|^2|\beta|^2$,
eine Größe auf der Achse der Mode, nicht zwei getrennte Ladungsteile an zwei Orten (Dok.~230) \BM.
Ein Zustand $\alpha|0\rangle+\beta|1\rangle$ ist in dieser Lesart eine einzige Konfiguration mit
einem Punkt $(z,r,\theta)$, nicht eine Hälfte hier und eine Hälfte dort. Die erste Annahme des
Einwands trifft auf diese Dichte nicht zu \SM.

\subsection{Eine Mode, eine Quelle}

Eine Mode hat eine einzige Zeitskala $\tilde T=1/m$ (Dok.~334); Überlagerungen verschiedener Moden
führen zu unverträglichen Zeitskalen, Überlagerung innerhalb einer Mode ist eine Konfiguration dieser
einen Mode (Dok.~333, 334). Daraus folgt in der Lesart der FFGFT: Eine Mode ist eine Quelle, nicht
ein Kontinuum von Quellen. Wechselwirkung koppelt verschiedene Moden, nicht eine Mode mit sich
selbst \SM.

In der Feldtheorie hat genau dieser Satz eine exakte Form \EM. Die Ladungsdichte ist ein Operator
$\hat\rho(x)=\hat\psi^\dagger(x)\hat\psi(x)$, und ihr Erwartungswert im Einteilchenzustand ist
$|\psi(x)|^2$; die Dichte ist also real im Sinn eines Erwartungswerts am Einzelsystem, so wie die
geschützte Messung sie zeigt. Der Coulomb-Term der Feldtheorie ist aber normalgeordnet,
\begin{equation}
  \hat H_C=\tfrac12\!\int\!\!\int \frac{k\,{:}\hat\rho(x)\hat\rho(x'){:}}{|x-x'|}\,d^3x\,d^3x' ,
\end{equation}
und ${:}\hat\rho\hat\rho{:}$ enthält zwei Vernichter. Auf jeden Einteilchenzustand angewandt ergibt er
null, für jede Überlagerung \BM. Das Prüfskript zeigt das auf einem Gitter mit exakter Rechnung im
Fock-Raum: Im Einteilchensektor ist der Erwartungswert des Zweikörperterms für alle Basiszustände
und für beliebige Überlagerungen null, während eine klassische Wolke derselben Dichte eine positive
Selbstenergie hätte; im Zweiteilchensektor ist die Wechselwirkung vorhanden \KM. Reale Dichte und
fehlende Selbstwechselwirkung schließen einander also nicht aus. Sie schließen einander nur aus,
wenn man die Dichte als klassisches Kontinuum liest.

Die bekannte Selbstenergie des Elektrons in der Quantenelektrodynamik ist etwas anderes: Sie stammt
aus der Kopplung an das Strahlungsfeld, also an andere Moden, und geht in die Massenrenormierung und
den Lamb-Shift ein \EM. Sie ist keine Abstoßung der Wolke gegen sich selbst und passt zur Lesart
\enquote{Wechselwirkung koppelt verschiedene Moden}.

\subsection{Was die FFGFT hier noch schuldet}

Dass die normalgeordnete Form die richtige ist, ist in der Feldtheorie gesetzt, nicht hergeleitet;
sie folgt aus der Forderung, dass das Vakuum keine Energie aus der eigenen Wechselwirkung trägt. Die
FFGFT liest dieselbe Struktur als Folge davon, dass eine Mode eine Quelle mit einer Zeitskala ist.
Eine Herleitung dieser Eigenschaft aus der Struktur des Trägers auf $T^4$ liegt nicht vor \SM.

Offen ist auch der Vielteilchenfall. Für $N$ Teilchen lebt $|\psi|^2$ im $3N$-dimensionalen
Konfigurationsraum. Die FFGFT führt Vielteilchenzustände im Produktraum der Moden (Dok.~230); dass die
Energiedichte dort dieselbe Lesart trägt wie für eine Mode, ist angenommen, nicht gezeigt \SM.

\subsection{Der gravitative Prüfstein}

Für die Gravitation folgt aus der Lesart eine Erwartung, die sich messen lässt. Wenn das Feld der
realisierten Lage folgt und eine Mode keine Quelle für sich selbst ist, gibt es den
Schrödinger-Newton-Term nicht. Die FFGFT erwartet daher in den vorgeschlagenen optomechanischen
Versuchen \cite{Yang2013,Grossardt2016} keine Verzerrung des Spektrums durch Selbstgravitation
\SM. Ein gemessener Schrödinger-Newton-Effekt würde dieser Lesart widersprechen. Die Erwartung ist
dieselbe wie die der gewöhnlichen Quantenmechanik ohne Selbstterm, aber sie ist begründet, nicht
nur übernommen: Sie folgt aus derselben Aussage über Quellen, die Dok.~020 für Page und Geilker
macht. Wie eine Masse vor jeder Wechselwirkung gravitativ wirkt, bleibt dabei offen (Dok.~020); die
Versuche zur gravitativ erzeugten Verschränkung prüfen diesen anderen Teil.

% ============================================================
\section{Gegenüberstellung}
% ============================================================
\begin{center}
\small
\begin{tabular}{@{}p{3.0cm}p{4.6cm}p{3.0cm}p{3.6cm}@{}}
\toprule
\textbf{Lesart} & \textbf{$|\psi|^2$ ist} & \textbf{Einzelablauf} & \textbf{Selbstwechselwirkung} \\
\midrule
Born & Wahrscheinlichkeitsdichte & zufällig & keine Frage \\
Schrödinger (wörtlich) & Ladungswolke im Ortsraum & stetig & vorhanden, \XM{} \\
Gao & effektive Dichte aus zufälliger Bewegung & zufällig, unstetig & keine \\
Bohm & Führungsfeld, Teilchen getrennt & deterministisch & keine \\
FFGFT & Energiedichte der inneren Konfiguration einer Mode & deterministisch & keine, eine Mode eine Quelle \SM \\
\bottomrule
\end{tabular}
\end{center}

Von den Lesarten, die $|\psi|^2$ eine reale Bedeutung am Einzelsystem geben, ist die FFGFT die einzige
in dieser Liste, die zugleich deterministisch ist und die Dichte selbst als Energie und nicht als
Führungsfeld oder Zeitmittel liest. Ihr Preis ist die offene Herleitung der Eigenschaft \enquote{eine
Mode, eine Quelle} aus dem Träger.

% ============================================================
\section{Übersicht}
% ============================================================
\begin{center}
\small
\begin{tabular}{@{}p{8.2cm}p{1.2cm}p{4.4cm}@{}}
\toprule
\textbf{Aussage} & \textbf{Status} & \textbf{Grund} \\
\midrule
Schrödinger: $-e|\psi|^2$ als Ladungsdichte & \QM & \cite{Schroedinger1926} \\
geschützte Messung liefert $|\psi|^2$ am Einzelsystem & \QM & \cite{AharonovVaidman1993,AAV1993} \\
Gaos Einwand: reale Dichte verlangt Selbstwechselwirkung & \QM & \cite{Gao2011} \\
Stärkemaße $8{,}5\times10^{-4}$ und $4{,}9\times10^{-89}$ & \KM & Prüfskript \\
$J_{1s}=5/8$\,Hartree, $J_{2s}=77/512$\,Hartree, 1S--2S verschoben um 12{,}9\,eV & \KM & Prüfskript \\
selbstabstoßende Ladungswolke im Ortsraum & \XM & 1S--2S, \cite{Parthey2011} \\
$r^2=4|\alpha|^2|\beta|^2$ auf der Achse einer Mode & \BM & Dok.~230 \\
${:}\hat\rho\hat\rho{:}$ verschwindet im Einteilchensektor & \BM & Prüfskript \\
FFGFT-Dichte trifft die erste Annahme des Einwands nicht & \SM & Dok.~020, 230 \\
eine Mode, eine Quelle; Herleitung aus dem Träger & \SM & offen \\
keine Schrödinger-Newton-Selbstgravitation erwartet & \SM & Prüfstein \\
\bottomrule
\end{tabular}
\end{center}

\begingroup\renewcommand{\chapter}[2]{}\small
\begin{thebibliography}{99}
\bibitem{Schroedinger1926} E.~Schrödinger, \textit{Quantisierung als Eigenwertproblem (Dritte und Vierte Mitteilung)}, Annalen der Physik \textbf{80} (1926) 437 und \textbf{81} (1926) 109.
\bibitem{BacciagaluppiValentini2009} G.~Bacciagaluppi, A.~Valentini, \textit{Quantum Theory at the Crossroads: Reconsidering the 1927 Solvay Conference}, Cambridge University Press 2009; arXiv:quant-ph/0609184.
\bibitem{AharonovVaidman1993} Y.~Aharonov, L.~Vaidman, \textit{Measurement of the Schrödinger wave of a single particle}, Physics Letters A \textbf{178} (1993) 38.
\bibitem{AAV1993} Y.~Aharonov, J.~Anandan, L.~Vaidman, \textit{Meaning of the wave function}, Physical Review A \textbf{47} (1993) 4616.
\bibitem{Gao2011} S.~Gao, \textit{The wave function and quantum reality}, arXiv:1108.1187 (2011).
\bibitem{Gao2017} S.~Gao, \textit{Is an electron a charge cloud? A reexamination of Schrödinger's charge density hypothesis}, Foundations of Science (2017), doi:10.1007/s10699-017-9521-3.
\bibitem{Salzman2005} P.~J.~Salzman, \textit{Investigation of the time dependent Schrödinger-Newton equation}, Dissertation, University of California at Davis, 2005.
\bibitem{Mead2000} C.~A.~Mead, \textit{Collective Electrodynamics: Quantum Foundations of Electromagnetism}, MIT Press 2000.
\bibitem{CramerMead2020} J.~G.~Cramer, C.~A.~Mead, \textit{Symmetry, transactions, and the mechanism of wave function collapse}, Symmetry \textbf{12} (2020) 1373; arXiv:2006.11365.
\bibitem{Parthey2011} C.~G.~Parthey et al., \textit{Improved measurement of the hydrogen 1S--2S transition frequency}, Physical Review Letters \textbf{107} (2011) 203001.
\bibitem{Diosi1984} L.~Diósi, \textit{Gravitation and quantum-mechanical localization of macro-objects}, Physics Letters A \textbf{105} (1984) 199.
\bibitem{Penrose1996} R.~Penrose, \textit{On gravity's role in quantum state reduction}, General Relativity and Gravitation \textbf{28} (1996) 581.
\bibitem{Yang2013} H.~Yang, H.~Miao, D.-S.~Lee, B.~Helou, Y.~Chen, \textit{Macroscopic quantum mechanics in a classical spacetime}, Physical Review Letters \textbf{110} (2013) 170401.
\bibitem{Grossardt2016} A.~Großardt, J.~Bateman, H.~Ulbricht, A.~Bassi, \textit{Optomechanical test of the Schrödinger-Newton equation}, Physical Review D \textbf{93} (2016) 096003.
\bibitem{PageGeilker1981} D.~N.~Page, C.~D.~Geilker, \textit{Indirect evidence for quantum gravity}, Physical Review Letters \textbf{47} (1981) 979.
\end{thebibliography}
\endgroup
